\documentclass[12pt]{article}
\usepackage[utf8]{inputenc}
\usepackage[T1]{fontenc}
\usepackage{hyperref}
\title{Ezaa Language Resources for K-12 Curriculum Development}
\author{Compiled by Claude Sonnet 5.5 (claude-sonnet-5-5) for Andr\'as Kornai}
\date{}
\begin{document}
\maketitle

\section{Introduction}
Ezaa (also Ezza, Eza; ISO 639-3 \texttt{eza}) is an Igboid language of Ebonyi State, Nigeria. It forms a dialect cluster with Izii, Ikwo and Mgbo, which are only marginally mutually intelligible with it. Speaker figures in the sources found conflict (roughly 590,000 for 2012 and about 1,590,000 as a 2020 projection). The search found only encyclopedic overview pages; no school materials in Ezaa were located.

\section{Primary Educational Resources}
No primers, textbooks or curricula specific to Ezaa were found in the search results.

\section{Linguistic and Reference Materials}
\begin{itemize}
\item Ezaa language (Wikipedia overview, classification and speaker estimates): \url{https://en.wikipedia.org/wiki/Ezaa_language}
\item Ezaa people (Wikipedia overview of the community in Ebonyi State): \url{https://en.wikipedia.org/wiki/Ezaa_people}
\end{itemize}

\section{Religious and Community Resources}
None found in the search results.

\section{Digital and Audio Resources}
None found in the search results.

\section{Curriculum Development Context}
Ezaa is closely related to standard Igbo, which is the likely written medium in schools of the area. Any Ezaa-medium curriculum would have to start from new documentation; the sources found are general overviews only and give conflicting claims about the alphabet, so orthography needs verification with the community.
\end{document}